\documentclass[12pt]{article}
\usepackage[margin=1in]{geometry}
\usepackage[T1]{fontenc}
\usepackage[utf8]{inputenc}
\usepackage{lmodern}
\usepackage{longtable}
\usepackage{booktabs}
\usepackage{array}
\usepackage{hyperref}
\usepackage{enumitem}
\usepackage{xcolor}

\hypersetup{
    colorlinks=true,
    linkcolor=blue,
    urlcolor=blue
}

\setlist[itemize]{leftmargin=1.5em}
\setlist[enumerate]{leftmargin=1.8em}

\title{HiPPO Results Report}
\author{}
\date{}

\begin{document}

\maketitle

\section*{What This Short Report Covers}

This version keeps the practical part of the work:

\begin{itemize}
    \item what datasets were run
    \item what results were obtained
    \item what the benchmark comparison showed
\end{itemize}

The long setup story was removed because it is not the main focus right now.

\section*{Datasets That Were Run}

During this work, three datasets were run:

\begin{enumerate}
    \item \texttt{mnist}
    \item \texttt{adding}
    \item \texttt{copying}
\end{enumerate}

If only two datasets are needed for the report, the best two to mention are:

\begin{itemize}
    \item \texttt{MNIST}
    \item \texttt{Copying}
\end{itemize}

This is because:

\begin{itemize}
    \item \texttt{MNIST} is the main real dataset that people know
    \item \texttt{Copying} is the dataset used for the 4-model comparison
\end{itemize}

\section*{What MNIST Is}

MNIST is a dataset of small black-and-white images of handwritten digits. The digits are 0, 1, 2, and so on up to 9.

The model looks at each image and tries to predict which digit it is.

In this project, \texttt{mnist} is the default dataset in the config, so when the project is run normally, it tries to use MNIST first.

\section*{MNIST Result}

After fixing the environment and manually providing the MNIST files, the project was run successfully on MNIST.

The result was:

\begin{itemize}
    \item training accuracy: \texttt{0.848}
    \item validation accuracy: \texttt{0.956}
    \item test accuracy: \texttt{0.961}
\end{itemize}

\section*{What The MNIST Result Means}

This means the model learned the task successfully.

The most important number is the test accuracy:

\begin{itemize}
    \item \texttt{0.961} means about \texttt{96.1\%} of the test images were classified correctly
\end{itemize}

In simple words:

\begin{itemize}
    \item the model was able to recognize handwritten digits well
    \item the training pipeline worked
    \item the GPU setup also worked
\end{itemize}

So MNIST is the clearest example showing that the project can run and produce a meaningful result.

\section*{What Adding Is}

\texttt{Adding} is not downloaded from the internet. It is an artificial dataset generated by the code itself.

The task is roughly this:

\begin{itemize}
    \item the model sees a sequence of numbers
    \item it must learn which values matter
    \item then it must produce the correct sum
\end{itemize}

This dataset is used to test whether the model can handle long sequences and keep useful information.

\section*{Adding Result}

The \texttt{adding} dataset was run successfully.

The visible final result from the run was:

\begin{itemize}
    \item loss: about \texttt{0.171}
\end{itemize}

\section*{What The Adding Result Means}

For this dataset, the main value is loss, not accuracy.

In simple words:

\begin{itemize}
    \item lower loss is better
    \item a loss of about \texttt{0.171} shows the model was learning the task
\end{itemize}

So the \texttt{adding} run shows that the project also worked on a synthetic sequence dataset, not only on image data.

\section*{What Copying Is}

\texttt{Copying} is also an artificial dataset generated inside the project.

This task checks sequence memory.

In simple words, the model must:

\begin{itemize}
    \item read a sequence
    \item remember important information from earlier in the sequence
    \item reproduce that information later
\end{itemize}

This is useful because it tests how well different model types handle memory over time.

\section*{Sequence Model Comparison Setup}

After getting the project to run, a comparison setup was added so different sequence-model families could be tested in the same training pipeline.

The following model families were used:

\begin{itemize}
    \item \texttt{gru} for \textbf{recurrence}
    \item \texttt{conv1d} for \textbf{convolution}
    \item \texttt{attention} for \textbf{attention}
    \item \texttt{legs} for \textbf{state-space / HiPPO}
\end{itemize}

Important note:

This repository does \textbf{not} contain a true S4 implementation. The closest built-in state-space model available here is \texttt{legs}, which is HiPPO-based. So in this project, the state-space comparison was done with \texttt{legs}, not with a full S4 model.

\section*{Benchmark Script}

To avoid running every benchmark manually, a script was added:

\begin{itemize}
    \item \texttt{run\_copying\_benchmarks.py}
\end{itemize}

This script:

\begin{enumerate}
    \item runs all four models on the \texttt{copying} dataset
    \item keeps the training settings consistent
    \item saves the full logs
    \item extracts runtime and final metrics into result files
\end{enumerate}

Generated output files:

\begin{itemize}
    \item \texttt{benchmark\_outputs/copying\_benchmark\_results.json}
    \item \texttt{benchmark\_outputs/copying\_benchmark\_results.md}
\end{itemize}

\section*{Benchmark Dataset Choice}

The \texttt{copying} dataset was used for the architecture comparison because:

\begin{itemize}
    \item it is built into the repository
    \item it does not need internet download
    \item it is a sequence task
    \item it is short enough that the attention model can run on the available GPU
\end{itemize}

The training command pattern used by the benchmark script was based on:

\begin{itemize}
    \item \texttt{dataset=copying}
    \item \texttt{dataset.num\_workers=0}
    \item \texttt{train.epochs=2}
\end{itemize}

The \texttt{dataset.num\_workers=0} setting was kept because this machine is Windows-based and the project otherwise runs into multiprocessing problems.

\section*{Copying Benchmark Results}

The following results were collected from the automated benchmark run:

\begin{center}
\begin{tabular}{l l r r r r}
\toprule
Model & Family & Runtime (s) & Train Acc & Val Acc & Test Acc \\
\midrule
\texttt{gru} & Recurrence & 485.03 & 0.109 & 0.009 & 0.008 \\
\texttt{conv1d} & Convolution & 28.67 & 0.104 & 0.101 & 0.099 \\
\texttt{attention} & Attention & 44.63 & 0.126 & 0.123 & 0.125 \\
\texttt{legs} & State-space / HiPPO & 639.89 & 0.094 & 0.132 & 0.130 \\
\bottomrule
\end{tabular}
\end{center}

\section*{What These Benchmark Results Mean}

These numbers compare two things:

\begin{itemize}
    \item how long the model took to run
    \item how well the model performed on the \texttt{copying} task
\end{itemize}

\subsection*{Runtime}

Runtime means training time.

From the results:

\begin{itemize}
    \item \texttt{conv1d} was the fastest: \texttt{28.67} seconds
    \item \texttt{attention} was also fairly fast: \texttt{44.63} seconds
    \item \texttt{gru} was much slower: \texttt{485.03} seconds
    \item \texttt{legs} was the slowest: \texttt{639.89} seconds
\end{itemize}

In simple words:

\begin{itemize}
    \item convolution was the fastest model here
    \item attention was also efficient
    \item recurrence and HiPPO took much longer in this experiment
\end{itemize}

\subsection*{Accuracy}

Accuracy tells us how often the model gave the correct answer.

The most useful value here is test accuracy:

\begin{itemize}
    \item \texttt{gru}: \texttt{0.008}
    \item \texttt{conv1d}: \texttt{0.099}
    \item \texttt{attention}: \texttt{0.125}
    \item \texttt{legs}: \texttt{0.130}
\end{itemize}

In simple words:

\begin{itemize}
    \item \texttt{legs} gave the best test result in this benchmark
    \item \texttt{attention} was close behind
    \item \texttt{conv1d} was lower
    \item \texttt{gru} performed poorly in this short run
\end{itemize}

\section*{Important Note About These Comparison Results}

These benchmark results should be explained carefully.

They do \textbf{not} mean:

\begin{itemize}
    \item convolution is always best
    \item recurrence is always bad
    \item attention is always best
    \item HiPPO is always best
\end{itemize}

Why not:

\begin{itemize}
    \item all models were run only for \texttt{2} epochs
    \item this was a short comparison
    \item no deep tuning was done for each model
\end{itemize}

So the fair conclusion is:

\begin{quote}
In this benchmark, convolution was the fastest model, while the HiPPO-based \texttt{legs} model achieved the best test accuracy on the \texttt{copying} task.
\end{quote}

\section*{Final Overall Conclusion}

The most important practical results from this work are:

\begin{enumerate}
    \item The project was run successfully on \texttt{MNIST}
    \item The project was also run on synthetic sequence datasets such as \texttt{adding} and \texttt{copying}
    \item A 4-model comparison was completed on the \texttt{copying} dataset
\end{enumerate}

The clearest result examples are:

\begin{itemize}
    \item \texttt{MNIST} test accuracy: \texttt{0.961}
    \item \texttt{Adding} loss: about \texttt{0.171}
    \item \texttt{Copying} benchmark: \texttt{conv1d} fastest, \texttt{legs} best test accuracy
\end{itemize}

So if the report needs a simple summary, it can say:

\begin{quote}
The project was successfully run on MNIST and on sequence datasets generated by the code. MNIST gave a strong result of about 96.1\% test accuracy. On the Copying benchmark, four model families were compared, and the fastest model was convolution, while the best test accuracy was achieved by the HiPPO-based \texttt{legs} model.
\end{quote}

\end{document}
